\chapter{Communication Complexity — Finite Probability Space}
%%% generated by the blueprint pipeline from TCSlib.CommunicationComplexity.NewmanTheorem.FiniteProbabilitySpace
\section{Overview}

This module develops the foundational infrastructure for finite probability spaces used
in the proof of Newman's theorem. It introduces the typeclasses
\texttt{CommunicationComplexity.FiniteMeasureSpace} and
\texttt{CommunicationComplexity.FiniteProbabilitySpace}, and establishes a library of
measure-real and integral lemmas (singleton decompositions, Cauchy--Schwarz, Markov's
inequality, product-space factorisation) that are used throughout the complexity-theoretic
arguments.

%%%%%%%%%%%%%%%%%%%%%%%%%%%%%%%%%%%%%%%%%%%%%%%%%%%%%%%%%%%%%%%%%%%%%%%%%%%%%%%
\section{Declarations}

\begin{definition}[Finite measurable space]
\lean{CommunicationComplexity.FiniteMeasureSpace}
\label{CommunicationComplexity.FiniteMeasureSpace}
\leanok
A typeclass for a measurable space $\Omega$ that is simultaneously finite (carries a
\texttt{Fintype} instance) and discrete (every subset is measurable).  It deliberately
does not bundle a measure, so that a single type can be equipped with many different
measures.
\end{definition}

\begin{definition}[Constructor for \texttt{FiniteMeasureSpace}]
\lean{CommunicationComplexity.FiniteMeasureSpace.of}
\label{CommunicationComplexity.FiniteMeasureSpace.of}
\leanok
\uses{CommunicationComplexity.FiniteMeasureSpace}
Given a type $\Omega$ that already carries \texttt{Fintype} and
\texttt{DiscreteMeasurableSpace} instances, \texttt{CommunicationComplexity.FiniteMeasureSpace.of} bundles them
into a \texttt{FiniteMeasureSpace} record without requiring any additional data.
\end{definition}

\begin{definition}[Finite probability space]
\lean{CommunicationComplexity.FiniteProbabilitySpace}
\label{CommunicationComplexity.FiniteProbabilitySpace}
\leanok
\uses{CommunicationComplexity.FiniteMeasureSpace}
A typeclass bundling a \texttt{MeasureSpace} structure on $\Omega$ together with a
\texttt{CommunicationComplexity.FiniteMeasureSpace} witness and an \texttt{IsProbabilityMeasure} instance for the
canonical volume measure.  This makes $\Omega$ simultaneously a finite discrete
measurable space and a probability space.
\end{definition}

\begin{definition}[Constructor from existing instances]
\lean{CommunicationComplexity.FiniteProbabilitySpace.of}
\label{CommunicationComplexity.FiniteProbabilitySpace.of}
\leanok
\uses{CommunicationComplexity.FiniteMeasureSpace.of, CommunicationComplexity.FiniteProbabilitySpace}
Given a type $\Omega$ already equipped with a \texttt{MeasureSpace}, \texttt{Fintype},
\texttt{DiscreteMeasurableSpace}, and \texttt{IsProbabilityMeasure} instance for
\texttt{volume}, this helper packages all of those into a \texttt{CommunicationComplexity.FiniteProbabilitySpace}
record.
\end{definition}

\begin{definition}[Constructor from an explicit measure]
\lean{CommunicationComplexity.FiniteProbabilitySpace.ofMeasure}
\label{CommunicationComplexity.FiniteProbabilitySpace.ofMeasure}
\leanok
\uses{CommunicationComplexity.FiniteMeasureSpace, CommunicationComplexity.FiniteProbabilitySpace}
Given a finite measurable space $\Omega$ and an unbundled probability measure $\mu$ on
it, \texttt{FiniteProbabilitySpace.ofMeasure} promotes $\mu$ to the canonical volume and
returns a \texttt{FiniteProbabilitySpace} instance.
\end{definition}

\begin{theorem}[Singleton decomposition of a finite measure]
\lean{CommunicationComplexity.FiniteMeasureSpace.measureReal_eq_sum_singletons}
\label{CommunicationComplexity.FiniteMeasureSpace.measureReal_eq_sum_singletons}
\leanok
\uses{CommunicationComplexity.FiniteMeasureSpace}
\difficulty{3}
Let $\Omega$ be a finite measurable space — finite and discrete, so that every subset is
measurable — and let $\mu$ be a finite measure on $\Omega$, writing $\mu^{\bbr}$ for its
real-valued mass function. Then for every subset $S \subseteq \Omega$,
\[
  \mu^{\bbr}(S) \;=\; \sum_{\omega \in \Omega}
\begin{cases} \mu^{\bbr}(\{\omega\}) & \text{if } \omega \in S, \\ 0 & \text{otherwise.}
\end{cases}
\]
That is, the measure of $S$ is the sum of the masses of its individual points.
\end{theorem}

\begin{theorem}[Measure of a preimage as a sum over fibers]
\lean{CommunicationComplexity.FiniteMeasureSpace.measureReal_preimage_eq_sum_fibers}
\label{CommunicationComplexity.FiniteMeasureSpace.measureReal_preimage_eq_sum_fibers}
\leanok
\uses{CommunicationComplexity.FiniteMeasureSpace, CommunicationComplexity.FiniteMeasureSpace.measureReal_eq_sum_singletons}
\difficulty{4}
Let $\Omega$ be a finite discrete measurable space, $\alpha$ a finite type, $\mu$ a
finite measure on $\Omega$, $Z : \Omega \to \alpha$ a function, and $P$ a predicate on
$\alpha$. Writing $\mu^{\bbr}$ for the real-valued measure obtained from $\mu$, the
measure of the set of points whose image under $Z$ satisfies $P$ decomposes as a sum
over the fibers of $Z$:
\[
  \mu^{\bbr}\{\omega \mid P(Z(\omega))\}
\;=\; \sum_{z \in \alpha} \begin{cases} \mu^{\bbr}\bigl(Z^{-1}\{z\}\bigr) & \text{if }
P(z), \\ 0 & \text{otherwise.} \end{cases}
\]
\end{theorem}

\begin{theorem}[Absolute continuity via singletons]
\lean{CommunicationComplexity.FiniteMeasureSpace.absolutelyContinuous_iff_forall_singletons}
\label{CommunicationComplexity.FiniteMeasureSpace.absolutelyContinuous_iff_forall_singletons}
\leanok
\uses{CommunicationComplexity.FiniteMeasureSpace}
\difficulty{4}
Let $\Omega$ be a finite measurable space, and let $\mu$ and $\nu$ be two measures on
$\Omega$. Then $\mu$ is absolutely continuous with respect to $\nu$, written $\mu \ll
\nu$, if and only if for every point $\omega \in \Omega$ with $\nu(\{\omega\}) = 0$ one
also has $\mu(\{\omega\}) = 0$.
\end{theorem}

\begin{theorem}[Square of the expectation is at most the expectation of the square]
\lean{CommunicationComplexity.FiniteMeasureSpace.sq_integral_le_integral_sq}
\label{CommunicationComplexity.FiniteMeasureSpace.sq_integral_le_integral_sq}
\leanok
\uses{CommunicationComplexity.FiniteMeasureSpace}
\difficulty{3}
Let $\Omega$ be a finite measurable space, let $\mu$ be a probability measure on
$\Omega$, and let $f : \Omega \to \bbr$ be any real-valued function. Then
\[
  \Bigl(\int_\Omega f\,d\mu\Bigr)^{2} \;\le\; \int_\Omega f^{2}\,d\mu.
\]
\end{theorem}

\begin{theorem}[Integral of a discrete-valued function as a sum over fibers]
\lean{CommunicationComplexity.FiniteMeasureSpace.integral_comp_eq_sum_measureReal_fibers}
\label{CommunicationComplexity.FiniteMeasureSpace.integral_comp_eq_sum_measureReal_fibers}
\leanok
\uses{CommunicationComplexity.FiniteMeasureSpace}
\difficulty{3}
Let $\Omega$ be a finite measurable space, $\alpha$ a finite type on which every subset
is measurable, $\mu$ a finite measure on $\Omega$, and $Z : \Omega \to \alpha$ and $f :
\alpha \to \bbr$ any functions. Then
\[
  \int_\Omega f(Z(\omega))\,d\mu(\omega)
  \;=\; \sum_{z \in \alpha} \mu\bigl(Z^{-1}(\{z\})\bigr)\, f(z),
\]
where $\mu(Z^{-1}(\{z\}))$ denotes the real number giving the measure of the fiber
$Z^{-1}(\{z\}) = \{\omega \in \Omega : Z(\omega) = z\}$.
\end{theorem}

\begin{theorem}[Law of total probability over the fibers of a random variable]
\lean{CommunicationComplexity.FiniteMeasureSpace.measureReal_eq_sum_cond_fiber_real}
\label{CommunicationComplexity.FiniteMeasureSpace.measureReal_eq_sum_cond_fiber_real}
\leanok
\uses{CommunicationComplexity.FiniteMeasureSpace}
\difficulty{4}
Let $\Omega$ be a finite measurable space equipped with a finite measure $\mu$, let
$\alpha$ be a finite set, let $Z : \Omega \to \alpha$ be a measurable ($\alpha$-valued)
map, and let $S \subseteq \Omega$. Writing $\mu^{\bbr}$ for the real-valued measure and
$\mu[\,\cdot \mid A\,]$ for the measure conditioned on an event $A$, we have
\[
  \mu^{\bbr}(S)
  \;=\; \sum_{z \in \alpha} \mu^{\bbr}\bigl(Z^{-1}(\{z\})\bigr)\,
        \cdot\, \mu\bigl[\,\cdot \mid Z^{-1}(\{z\})\,\bigr]^{\bbr}(S),
\]
so the total mass of $S$ is recovered by summing, over the values $z$ of $Z$, the mass
of the fiber $Z^{-1}(\{z\})$ times the conditional mass of $S$ given that fiber.
\end{theorem}

\begin{theorem}[Uniform measure of a set equals its relative cardinality]
\lean{CommunicationComplexity.uniformOn_univ_measureReal_eq_card_filter}
\label{CommunicationComplexity.uniformOn_univ_measureReal_eq_card_filter}
\leanok
\difficulty{2}
Let $\Omega$ be a nonempty finite set, equipped with the uniform probability measure
that assigns each point mass $1/\abs{\Omega}$, and let $S \subseteq \Omega$. Then the
measure of $S$, as a real number, is
\[
  \frac{\abs{\{\omega \in \Omega \mid \omega \in S\}}}{\abs{\Omega}},
\]
the number of elements of $\Omega$ lying in $S$ divided by the total number of elements
of $\Omega$.
\end{theorem}

\begin{theorem}[Uniform measure of a set as a cardinality ratio]
\lean{CommunicationComplexity.uniformOn_univ_measureReal_eq_card_subtype}
\label{CommunicationComplexity.uniformOn_univ_measureReal_eq_card_subtype}
\leanok
\uses{CommunicationComplexity.uniformOn_univ_measureReal_eq_card_filter}
\difficulty{2}
Let $\Omega$ be a nonempty finite set, equipped with its discrete measurable structure,
and let $S \subseteq \Omega$. Under the uniform probability measure on $\Omega$, the
measure of $S$, regarded as a real number, equals
\[
\frac{\abs{S}}{\abs{\Omega}},
\]
where $\abs{S}$ is the number of elements of $\Omega$ lying in $S$ and $\abs{\Omega}$ is
the total number of elements of $\Omega$.
\end{theorem}

\begin{theorem}[Nonemptiness of a finite probability space]
\lean{CommunicationComplexity.FiniteProbabilitySpace.nonempty}
\label{CommunicationComplexity.FiniteProbabilitySpace.nonempty}
\leanok
\uses{CommunicationComplexity.FiniteProbabilitySpace}
\difficulty{1}
If a type $\Omega$ carries the structure of a finite probability space — that is, it is
finite and discrete, and is equipped with a probability measure — then $\Omega$ is
nonempty.
\end{theorem}

\begin{definition}[Probability mass function of a finite probability space]
\lean{CommunicationComplexity.FiniteProbabilitySpace.toPMF}
\label{CommunicationComplexity.FiniteProbabilitySpace.toPMF}
\leanok
\uses{CommunicationComplexity.FiniteProbabilitySpace}
For a finite probability space $\Omega$, \texttt{toPMF} converts the canonical volume
measure to a probability mass function $\mathrm{PMF}\,\Omega$, assigning to each
$\omega$ the real weight $\mu(\{\omega\})$.
\end{definition}

\begin{theorem}[Measure of a set as a sum of point masses]
\lean{CommunicationComplexity.FiniteProbabilitySpace.measure_eq}
\label{CommunicationComplexity.FiniteProbabilitySpace.measure_eq}
\leanok
\uses{CommunicationComplexity.FiniteProbabilitySpace, CommunicationComplexity.FiniteProbabilitySpace.toPMF}
\difficulty{3}
Let $\Omega$ be a finite probability space, write $\mu$ for its canonical probability
measure, and let $p(\omega) = \mu(\{\omega\})$ denote the associated probability mass
function. Then for every subset $S \subseteq \Omega$,
\[
\mu(S) = \sum_{\omega \in S} p(\omega).
\]
\end{theorem}

\begin{theorem}[Singleton masses of a finite probability space sum to one]
\lean{CommunicationComplexity.FiniteProbabilitySpace.hasSum_measure_singletons}
\label{CommunicationComplexity.FiniteProbabilitySpace.hasSum_measure_singletons}
\leanok
\uses{CommunicationComplexity.FiniteProbabilitySpace}
\difficulty{3}
Let $\Omega$ be a finite probability space, so that $\Omega$ is finite and discrete and
its canonical volume measure $\mathbb{P}$ is a probability measure, and let $\alpha$ be
a finite type equipped with a bijection $e : \Omega \xrightarrow{\sim} \alpha$. Then the
family of singleton masses $\bigl(\mathbb{P}(\{e^{-1}(a)\})\bigr)_{a \in \alpha}$ is
summable with sum $1$.
\end{theorem}

\begin{theorem}[Probability mass function of a product factors]
\lean{CommunicationComplexity.FiniteProbabilitySpace.pmf_prod}
\label{CommunicationComplexity.FiniteProbabilitySpace.pmf_prod}
\leanok
\uses{CommunicationComplexity.FiniteProbabilitySpace, CommunicationComplexity.FiniteProbabilitySpace.toPMF}
\difficulty{3}
Let $\Omega_1$ and $\Omega_2$ be finite probability spaces, and give the product
$\Omega_1 \times \Omega_2$ its induced finite probability space structure. Then for
every $x \in \Omega_1$ and $y \in \Omega_2$, the probability mass function of the
product assigns to the pair $(x, y)$ the product of the individual masses; that is, its
value at $(x, y)$ equals the value of the mass function of $\Omega_1$ at $x$ times the
value of the mass function of $\Omega_2$ at $y$.
\end{theorem}

\begin{theorem}[Product measure of a rectangle]
\lean{CommunicationComplexity.FiniteProbabilitySpace.measureReal_prod}
\label{CommunicationComplexity.FiniteProbabilitySpace.measureReal_prod}
\leanok
\uses{CommunicationComplexity.FiniteProbabilitySpace}
\difficulty{2}
Let $\Omega_1$ and $\Omega_2$ be finite probability spaces, with respective probability
measures $\bbP_1$ and $\bbP_2$, and let $A \subseteq \Omega_1$ and $B \subseteq
\Omega_2$ be any subsets. Then, under the product probability measure $\bbP$ on
$\Omega_1 \times \Omega_2$, the measure of the rectangle $A \times B$ factors as
\[
  \bbP(A \times B) \;=\; \bbP_1(A) \cdot \bbP_2(B),
\]
where each quantity is taken as a real number.
\end{theorem}

\begin{theorem}[Finite additivity of the probability measure]
\lean{CommunicationComplexity.FiniteProbabilitySpace.measureReal_iUnion_fintype}
\label{CommunicationComplexity.FiniteProbabilitySpace.measureReal_iUnion_fintype}
\leanok
\uses{CommunicationComplexity.FiniteProbabilitySpace}
\difficulty{2}
Let $\Omega$ be a finite probability space and let $\iota$ be a finite index set. If
$(A_i)_{i \in \iota}$ is a pairwise disjoint family of subsets of $\Omega$, then the
probability of their union is the sum of their probabilities:
\[
\mathbb{P}\!\Bigl(\bigcup_{i \in \iota} A_i\Bigr) \;=\; \sum_{i \in \iota}
\mathbb{P}(A_i),
\]
where $\mathbb{P}(\,\cdot\,) \in \bbr$ denotes the real-valued probability measure of
$\Omega$.
\end{theorem}

\begin{theorem}[Measure of a preimage of a finite set]
\lean{CommunicationComplexity.FiniteProbabilitySpace.measureReal_preimage_finset}
\label{CommunicationComplexity.FiniteProbabilitySpace.measureReal_preimage_finset}
\leanok
\uses{CommunicationComplexity.FiniteProbabilitySpace}
\difficulty{3}
Let $\Xi$ be a finite probability space and let $\Omega$ be a measurable space in which
every subset is measurable. For any map $\varphi \colon \Xi \to \Omega$ and any finite
set $T \subseteq \Omega$, the probability that $\varphi$ lands in $T$ decomposes as a
sum over the points of $T$:
\[
  \mathrm{volume}^{\mathbb{R}}\bigl(\varphi^{-1}(T)\bigr)
  \;=\; \sum_{a \in T} \mathrm{volume}^{\mathbb{R}}\bigl(\varphi^{-1}(\{a\})\bigr),
\]
where $\mathrm{volume}^{\mathbb{R}}$ denotes the real-valued measure of the underlying
probability measure on $\Xi$.
\end{theorem}

\begin{theorem}[Additivity of the measure over a finite set]
\lean{CommunicationComplexity.FiniteProbabilitySpace.measureReal_finset}
\label{CommunicationComplexity.FiniteProbabilitySpace.measureReal_finset}
\leanok
\uses{CommunicationComplexity.FiniteProbabilitySpace, CommunicationComplexity.FiniteProbabilitySpace.measureReal_preimage_finset}
\difficulty{1}
Let $\Omega$ be a finite probability space, and let $T$ be a finite subset of $\Omega$.
Writing $\mathbb{P}$ for the real-valued probability measure on $\Omega$, the measure of
$T$ is the sum of the measures of its singletons:
\[
  \mathbb{P}(T) \;=\; \sum_{a \in T} \mathbb{P}(\{a\}).
\]
\end{theorem}

\begin{theorem}[Integral as a probability-weighted sum on a finite probability space]
\lean{CommunicationComplexity.FiniteProbabilitySpace.integral_eq_pmf_sum}
\label{CommunicationComplexity.FiniteProbabilitySpace.integral_eq_pmf_sum}
\leanok
\uses{CommunicationComplexity.FiniteProbabilitySpace, CommunicationComplexity.FiniteProbabilitySpace.toPMF}
\difficulty{1}
Let $\Omega$ be a finite probability space, and let $p(\omega)$ denote the probability
mass it assigns to each point $\omega \in \Omega$. Then for every function $f \colon
\Omega \to \bbr$,
\[
  \int_\Omega f(\omega)\,d\omega \;=\; \sum_{\omega \in \Omega} p(\omega)\, f(\omega),
\]
so that the expectation of $f$ under the canonical probability measure equals the finite
$p$-weighted sum of its values.
\end{theorem}

\begin{theorem}[Squared mean at most mean of square on a finite probability space]
\lean{CommunicationComplexity.FiniteProbabilitySpace.sq_integral_le_integral_sq}
\label{CommunicationComplexity.FiniteProbabilitySpace.sq_integral_le_integral_sq}
\leanok
\uses{CommunicationComplexity.FiniteMeasureSpace.sq_integral_le_integral_sq, CommunicationComplexity.FiniteProbabilitySpace}
\difficulty{1}
Let $\Omega$ be a finite probability space, and let $f\colon\Omega\to\bbr$ be any
real-valued function. Then the square of the expectation of $f$ is at most the
expectation of its square:
\[
  \Bigl(\int_\Omega f\,\Bigr)^2 \;\le\; \int_\Omega f^2,
\]
where both integrals are taken with respect to the underlying probability measure.
\end{theorem}

\begin{theorem}[Marginal of a coordinate integral over a product space]
\lean{CommunicationComplexity.FiniteProbabilitySpace.integral_comp_eval}
\label{CommunicationComplexity.FiniteProbabilitySpace.integral_comp_eval}
\leanok
\uses{CommunicationComplexity.FiniteProbabilitySpace}
\difficulty{3}
Let $\iota$ be a finite index type and let $\Omega$ be a finite probability space, so
that the product $\Omega^{\iota}$ of copies of $\Omega$ indexed by $\iota$ carries the
corresponding product probability measure. Fix an index $i \in \iota$ and a function $f
: \Omega \to \bbr$. Then integrating the map $\omega_{\bullet} \mapsto f(\omega_i)$
against the product measure on $\Omega^{\iota}$ gives the same value as integrating $f$
against the measure on $\Omega$:
\[
\int_{\Omega^{\iota}} f(\omega_i) \, d\omega_{\bullet} \;=\; \int_{\Omega} f(\omega)\,
d\omega.
\]
\end{theorem}

\begin{theorem}[Product measure of a box factorises]
\lean{CommunicationComplexity.FiniteProbabilitySpace.measureReal_pi_univ}
\label{CommunicationComplexity.FiniteProbabilitySpace.measureReal_pi_univ}
\leanok
\uses{CommunicationComplexity.FiniteProbabilitySpace}
\difficulty{2}
Let $\iota$ be a finite index set, and for each $i \in \iota$ let $\Omega_i$ be a finite
probability space carrying its canonical probability measure. For any choice of subsets
$s_i \subseteq \Omega_i$, the product measure on the finite product $\prod_{i \in \iota}
\Omega_i$ assigns to the box $\prod_{i \in \iota} s_i$ the value
\[
\mathbb{P}\!\Bigl(\prod_{i \in \iota} s_i\Bigr) \;=\; \prod_{i \in \iota}
\mathbb{P}(s_i),
\]
where both sides are read as real numbers.
\end{theorem}

\begin{theorem}[Measure of a set as the integral of its indicator]
\lean{CommunicationComplexity.FiniteProbabilitySpace.measureReal_eq_integral_indicator_one}
\label{CommunicationComplexity.FiniteProbabilitySpace.measureReal_eq_integral_indicator_one}
\leanok
\uses{CommunicationComplexity.FiniteProbabilitySpace}
\difficulty{1}
Let $\Omega$ be a finite probability space and let $S \subseteq \Omega$. Then the
(real-valued) probability of $S$ equals the integral over $\Omega$ of the indicator
function of $S$, taken with value $1$ on $S$ and $0$ elsewhere:
\[
\mathbb{P}(S) \;=\; \int_{\Omega} \mathbf{1}_S(\omega)\, d\omega,
\]
where $\mathbf{1}_S \colon \Omega \to \bbr$ is the function equal to $1$ at points of
$S$ and $0$ at points outside $S$.
\end{theorem}

\begin{theorem}[Probability as the integral of an indicator]
\lean{CommunicationComplexity.FiniteProbabilitySpace.measureReal_eq_integral_indicator}
\label{CommunicationComplexity.FiniteProbabilitySpace.measureReal_eq_integral_indicator}
\leanok
\uses{CommunicationComplexity.FiniteProbabilitySpace, CommunicationComplexity.FiniteProbabilitySpace.measureReal_eq_integral_indicator_one}
\difficulty{2}
Let $\Omega$ be a finite probability space and let $S \subseteq \Omega$ be any subset.
Then the probability of $S$ equals the integral of the indicator function of $S$ against
the probability measure:
\[
\mathbb{P}(S) \;=\; \int_\Omega \bigl(\text{if } \omega \in S \text{ then } 1 \text{
else } 0\bigr)\,d\omega,
\]
where the left-hand side is the real number obtained from the measure of $S$.
\end{theorem}

\begin{theorem}[Probability mass function sums to one]
\lean{CommunicationComplexity.FiniteProbabilitySpace.pmf_toReal_sum_eq_one}
\label{CommunicationComplexity.FiniteProbabilitySpace.pmf_toReal_sum_eq_one}
\leanok
\uses{CommunicationComplexity.FiniteProbabilitySpace, CommunicationComplexity.FiniteProbabilitySpace.toPMF}
\difficulty{3}
Let $\Omega$ be a finite probability space, and for each $\omega \in \Omega$ let
$p(\omega)$ denote the weight that the associated probability mass function assigns to
$\omega$, viewed as a real number. Then these weights sum to one: \[ \sum_{\omega \in
\Omega} p(\omega) = 1. \]
\end{theorem}

\begin{theorem}[Nonnegativity of probability mass weights]
\lean{CommunicationComplexity.FiniteProbabilitySpace.pmf_toReal_nonneg}
\label{CommunicationComplexity.FiniteProbabilitySpace.pmf_toReal_nonneg}
\leanok
\uses{CommunicationComplexity.FiniteProbabilitySpace, CommunicationComplexity.FiniteProbabilitySpace.toPMF}
\difficulty{1}
Let $\Omega$ be a finite probability space, and let $\mu(\{\omega\})$ denote the
probability mass that its associated probability mass function assigns to a point
$\omega \in \Omega$. Then this weight, taken as a real number, is nonnegative: \[ 0 \le
\mu(\{\omega\}). \]
\end{theorem}

\begin{theorem}[Existence of a point of positive mass]
\lean{CommunicationComplexity.FiniteProbabilitySpace.exists_pmf_toReal_pos}
\label{CommunicationComplexity.FiniteProbabilitySpace.exists_pmf_toReal_pos}
\leanok
\uses{CommunicationComplexity.FiniteProbabilitySpace, CommunicationComplexity.FiniteProbabilitySpace.pmf_toReal_sum_eq_one, CommunicationComplexity.FiniteProbabilitySpace.toPMF}
\difficulty{3}
Let $\Omega$ be a finite probability space, and let $\mu$ denote its associated
probability mass function, assigning to each point $\omega \in \Omega$ the real weight
$\mu(\{\omega\}) \ge 0$. Then there exists a point $\omega \in \Omega$ with
$\mu(\{\omega\}) > 0$.
\end{theorem}

\begin{theorem}[Pointwise upper bound descends to the integral]
\lean{CommunicationComplexity.FiniteProbabilitySpace.integral_le_of_le}
\label{CommunicationComplexity.FiniteProbabilitySpace.integral_le_of_le}
\leanok
\uses{CommunicationComplexity.FiniteProbabilitySpace, CommunicationComplexity.FiniteProbabilitySpace.integral_eq_pmf_sum, CommunicationComplexity.FiniteProbabilitySpace.pmf_toReal_nonneg, CommunicationComplexity.FiniteProbabilitySpace.pmf_toReal_sum_eq_one, CommunicationComplexity.FiniteProbabilitySpace.toPMF}
\difficulty{3}
Let $\Omega$ be a finite probability space, and let $f : \Omega \to \bbr$ be a function
satisfying $f(\omega) \le c$ for every $\omega \in \Omega$, where $c \in \bbr$ is a
constant. Then the integral of $f$ with respect to the probability measure on $\Omega$
satisfies $\int_\Omega f \le c$.
\end{theorem}

\begin{theorem}[Markov's inequality on a finite probability space]
\lean{CommunicationComplexity.FiniteProbabilitySpace.measureReal_ge_le_integral_div}
\label{CommunicationComplexity.FiniteProbabilitySpace.measureReal_ge_le_integral_div}
\leanok
\uses{CommunicationComplexity.FiniteProbabilitySpace}
\difficulty{3}
Let $\Omega$ be a finite probability space and let $f : \Omega \to \bbr$ be a
nonnegative function, so that $f(\omega) \ge 0$ for every $\omega \in \Omega$. Then for
every real $\varepsilon > 0$, the probability that $f$ attains at least $\varepsilon$ is
bounded by the mean of $f$ divided by $\varepsilon$:
\[
  \E\!\left[\,\{\omega \mid \varepsilon \le f(\omega)\}\,\right]
  \;\le\; \frac{\int_\Omega f(\omega)\,d\omega}{\varepsilon}.
\]
\end{theorem}

\begin{theorem}[Strict lower bound for an integral]
\lean{CommunicationComplexity.FiniteProbabilitySpace.lt_integral_of_lt}
\label{CommunicationComplexity.FiniteProbabilitySpace.lt_integral_of_lt}
\leanok
\uses{CommunicationComplexity.FiniteProbabilitySpace, CommunicationComplexity.FiniteProbabilitySpace.exists_pmf_toReal_pos, CommunicationComplexity.FiniteProbabilitySpace.integral_eq_pmf_sum, CommunicationComplexity.FiniteProbabilitySpace.pmf_toReal_nonneg, CommunicationComplexity.FiniteProbabilitySpace.pmf_toReal_sum_eq_one, CommunicationComplexity.FiniteProbabilitySpace.toPMF}
\difficulty{3}
Let $\Omega$ be a finite probability space and let $f : \Omega \to \bbr$ be a function.
If there is a constant $c \in \bbr$ with $c < f(\omega)$ for every $\omega \in \Omega$,
then $c < \int_\Omega f$, where the integral is taken with respect to the probability
measure on $\Omega$.
\end{theorem}
